%% ============================================================
\section{Trace Analysis}\label{sec:analysis}
%% ============================================================


% TODO: Formalise the synthesis encoding and discuss scalability.
As outlined in \Cref{sec:design:pipeline}, we position trace analysis as a synthesis problem: given the model and an observable trace, we ask a solver to find a sequence of internal actions that, interleaved with the observed externals, is admissible by the model. A satisfying solution is a witness execution containing the internal actions that explains the observation; an unsatisfiable query indicates that no model-admissible execution explains the trace, and is evidence that the model is incorrect or incomplete. We encoded the model as a set of constraints over a single linear trace that interleaves the actions of all cores, and require that the solver's solution contains the observed external actions in the same order as they appear in the observable trace.

For the purpose of our efforts, we used Rosette(a program synthesis library built using Racket) to encode the synthesis problem. The greatest challenge in this process was to ensure that the synthesis problem is scalable and does not explode computationally. We tried to achieve this by the same reduction arguments mentioned in section~\ref{sec:reduction}. This was not entirely successful and we were only able to synthesize simple traces with a small number of steps (1-2 internal actions) and upto 2 cores in certain cases. 
We believe that this can be improved in several complementary ways:

\begin{itemize}
  \item \textbf{Coarser walk structure and overlap-based reasoning.}
  Instead of letting the solver choose each fine-grained walk micro-step, one can introduce a single \emph{walk} macro-step that stands for an entire walk history, with internal ordering resolved only where it matters.
  With concurrent walks, many pairs are effectively \emph{quiescent}: walks that do not overlap a modification to the paging structures they consult can often be treated as atomic with respect to those updates, which shrinks the interleaving space.
  What then matters is (i) whether a walk is \emph{in flight} over a virtual-address range while paging structures for that range are being modified, and (ii) where in that window the update becomes architecturally visible relative to the walk.

  \item \textbf{Per-core abstraction and pruned action sets.}
  Another reduction is to ignore, where sound, non-write events from other cores when reasoning about a given core's translation story, or to separate each core's contribution and only merge at explicit synchronization points.
  More generally, the transition relation need not expose every action at every state: domain knowledge from the reduction argument (Section~\ref{sec:reduction}) can cut down the actions eligible from each state, rather than enumerating a full Cartesian product of interleavings.
\end{itemize}

